Somatic structural variation (SV) is a hallmark of cancer, typically defined as the insertion, deletion, or rearrangement of genomic segments longer than 50 bp~\cite{cosenza}. Somatic SVs vary in size and complexity far more than single-nucleotide variants and short indels, from a single deletion to catastrophic events that shuffle fragments of several chromosomes~\cite{chromothripsis}. Different cancer types exhibit characteristic patterns of them: 16 signatures of structural variation emerged across 2,658 cancers from 38 types~\cite{pcawg}. As a tumor grows, somatic events accumulate, so the tumor cells of one patient are not copies of each other~\cite{ith}: the junctions they carry are distinct from those confined to one region, metastasis, or time point, and that separation constitutes the case's subclonal architecture.

Most somatic variant callers still operate on short reads, which miss structural variants within repetitive regions, and recent benchmarks report low sensitivity and high false-positive rates~\cite{zook,wagner,parliament2}. Long-read technologies, such as Oxford Nanopore Technologies (ONT) and Pacific Biosciences (PacBio), span those repeats and address the limitations of short-read SV calling~\cite{sedlazeck,logsdon}. Read accuracy matters as much as read length here, because a method that compares samples by sequence identity rather than by mapped position requires a locus to spell the same way in every read set covering it. PacBio HiFi reads reach 99.8\% accuracy at about 13.5 kb~\cite{wenger}, and on the PacBio Revio data used here we measure a median read length of 16.2 kb at an overall error of 0.018\%, at which about 96\% of 200 bp stretches are error-free.

Still, long-read somatic SV callers rely on a reference genome to map both the tumor and the normal reads before deciding which sites are somatic~\cite{severus,savana,nanomonsv}. This introduces reference bias: reads from regions that diverge from the reference map less reliably, so SVs are undercalled exactly where the tumor departs from it most~\cite{refbias}. As a solution, the de Bruijn graph (dBG) avoids the reference entirely, representing a read set as a graph of overlapping $k$-mers~\cite{debruijn}, and to compare more than one genome, Cortex colors each node by the samples its $k$-mers came from~\cite{cortex}. Building the graph from $k$-mers, however, carries its own drawbacks: the resolution is fixed in advance by $k$, a single substitution replaces $k$ consecutive $k$-mers, and every position of every read contributes one.

Two families of methods address these drawbacks. Sketching keeps only a subset of the $k$-mers: minimizers retain the smallest $k$-mer of each window~\cite{roberts}, syncmers select by internal position and tolerate errors better~\cite{syncmers}, and de Bruijn graphs have been built in minimizer space on this idea~\cite{rustmdbg}. What sketching cannot promise is coverage. Two sequences are guaranteed to share a minimizer only when their exact match reaches $w+k-1$ bases, and no $k$-mer sketching scheme guarantees a fixed distance between consecutive selected substrings~\cite{lcpan}. Locally Consistent Parsing (LCP), which predates both minimizers and MinHash~\cite{sv94a,sv94b}, partitions the string instead of sampling it: cores are distributed uniformly over the input, identical cores carry identical labels, and every character lies in at least one core~\cite{lcpan}. Nothing is discarded, so no locus can fall between the selected pieces~\cite{gencore}. Core length follows the sequence and grows by a factor of about 2.34 per level~\cite{lcpan}, which turns resolution into a level rather than a guessed $k$; and sketch distances are known to blur precisely where genomes are most similar~\cite{gencore}, which is the regime of several tumors from one patient.

A further limitation is that these tools accept one tumor--normal pair at a time. Recovering the phylogeny of $N$ tumor lines means running them $N$ times, so each junction is discovered separately in every tumor that carries it, and whether two tumors share one is decided afterwards by intersecting coordinate intervals. A shared event placed differently in two runs is then reported as two events.

In this work, we present \tool{}, a somatic SV caller that builds a personalized pangenome from all tumor lines of one patient at once, using a colored graph whose nodes are LCP cores rather than $k$-mers. It therefore makes the somatic decision before any read is mapped: \tool{} takes one normal and $N$ tumor read sets, counts every core in each of them, and extracts the candidate somatic events of all lines in a single pass.

We show that \tool{} is comparable to the published calls of Severus, SAVANA and nanomonsv on the CASTLE panel while mapping 30--100 times less sequence, and that a case of twenty tumors is called faster and slightly more accurately in one graph than in twenty pairwise runs. Scored against that panel's own consensus, which eleven published callers built and \tool{} did not, it ranks third of twelve; rebuilt with \tool{} contributing on the same terms as the others, it ranks first. End-to-end, it is about 6 times faster than mapping both read sets to a reference, at 1.6 times the peak memory.

